% GENERATED FILE. Edit canonical lessons, then run npm run generate:books.
% canonical-source-hash: fnv1a64:1d14e3066fcf53a6
% canonical-lessons: AR-C307-jaysh, AR-C307-jundi, AR-C307-harb, AR-C307-maraka, AR-C307-haraba-war

\chapter{Armies and wars}
\label{ch:ar-armies-and-wars}
\hlchaptermodality{307}

\begin{chapteropening}
\textbf{By the end of this chapter:} \emph{I can say an army, a soldier, a war, a battle and to wage war.}

Complete the last lesson of chapter 307: I can say an army, a soldier, a war, a battle and to wage war.
\end{chapteropening}

\section[\texorpdfstring{\ar{جيش} (jaysh)}{جيش (jaysh)}]{\ar{جيش} (jaysh) — an army}

\label{lesson:AR-C307-jaysh}



\begin{quote}
Before the new one: say the Arabic for to escort a bride, then the Arabic for to ululate for joy.
\end{quote}

\subsection*{You'll want to know: \ar{جيش}}
\textbf{\ar{جيش}} — \emph{jaysh} — \textquotedblleft{}an army\textquotedblright{}.

\begin{quote}
\ar{عبر} \ar{الجيش} \ar{النهر} \ar{في} \ar{الليل}.
\end{quote}

\begin{quote}
\emph{ʿabara al-jaysh an-nahr fī al-layl.}
\end{quote}

\begin{quote}
The army crossed the river at night.
\end{quote}

Plural: \ar{جيوش}, \emph{juyūsh}.

\begin{cousinweb}
\textbf{\ar{جيش}} (\emph{jaysh}) is built on the native Arabic root \emph{j-y-sh} (Lane), where \emph{jāsha} means \textquotedblleft{}to rise, to heave, to surge\textquotedblright{}. Lane gives \emph{jaysh} as an army, or men going out to war.
\end{cousinweb}

\begin{grammarlens}[title={What you've built}]
One more word for telling what happened.
\end{grammarlens}

\subsection*{Your turn}
\begin{itemize}
\raggedright
  \item \emph{Say it:} \emph{jaysh}
  \item \emph{Say it:} \emph{jaysh}, once more
  \item \emph{Say it:} \emph{ʿabara al-jaysh an-nahr fī al-layl.} — The army crossed the river at night.
  \item \emph{Recall:} say the Arabic for to escort a bride, then the Arabic for to ululate for joy, then say \emph{jaysh} again
\end{itemize}

\begin{tcolorbox}[breakable,colback=teal!4,colframe=teal!35!black,title={Before you move on}]
What does \ar{جيش} mean? (\textquotedblleft{}An army\textquotedblright{}.) Say it once more.
\end{tcolorbox}
\section[\texorpdfstring{\ar{جندي} (jundī)}{جندي (jundī)}]{\ar{جندي} (jundī) — a soldier}

\label{lesson:AR-C307-jundi}



\begin{quote}
Before the new one: say the Arabic for to ululate for joy, then the Arabic for an army.
\end{quote}

\subsection*{You'll want to know: \ar{جندي}}
\textbf{\ar{جندي}} — \emph{jundī} — \textquotedblleft{}a soldier\textquotedblright{}.

\begin{quote}
\ar{كان} \ar{جدي} \ar{جنديا} \ar{في} \ar{شبابه}.
\end{quote}

\begin{quote}
\emph{kāna jaddī jundiyyan fī shabābihi.}
\end{quote}

\begin{quote}
My grandfather was a soldier in his youth.
\end{quote}

Plural: \ar{جنود}, \emph{junūd}.

\begin{cousinweb}
\textbf{\ar{جندي}} (\emph{jundī}) is built on \emph{jund}, \textquotedblleft{}troops, a host\textquotedblright{}, filed under the root \emph{j-n-d} (Lane), with the nisba ending \emph{-ī}: \textquotedblleft{}one of the troops\textquotedblright{}.
\end{cousinweb}

\begin{grammarlens}[title={What you've built}]
One more word for telling what happened.
\end{grammarlens}

\subsection*{Your turn}
\begin{itemize}
\raggedright
  \item \emph{Say it:} \emph{jundī}
  \item \emph{Say it:} \emph{jundī}, once more
  \item \emph{Say it:} \emph{kāna jaddī jundiyyan fī shabābihi.} — My grandfather was a soldier in his youth.
  \item \emph{Recall:} say the Arabic for to ululate for joy, then the Arabic for an army, then say \emph{jundī} again
\end{itemize}

\begin{tcolorbox}[breakable,colback=teal!4,colframe=teal!35!black,title={Before you move on}]
What does \ar{جندي} mean? (\textquotedblleft{}A soldier\textquotedblright{}.) Say it once more.
\end{tcolorbox}
\section[\texorpdfstring{\ar{حرب} (ḥarb)}{حرب (ḥarb)}]{\ar{حرب} (ḥarb) — a war}

\label{lesson:AR-C307-harb}



\begin{quote}
Before the new one: say the Arabic for an army, then the Arabic for a soldier.
\end{quote}

\subsection*{You'll want to know: \ar{حرب}}
\textbf{\ar{حرب}} — \emph{ḥarb} — \textquotedblleft{}a war\textquotedblright{}.

\begin{quote}
\ar{كان} \ar{جدي} \ar{صغيرا} \ar{في} \ar{زمن} \ar{الحرب}.
\end{quote}

\begin{quote}
\emph{kāna jaddī ṣaghīran fī zaman al-ḥarb.}
\end{quote}

\begin{quote}
My grandfather was a small boy in wartime.
\end{quote}

Plural: \ar{حروب}, \emph{ḥurūb}.

\ar{حرب} is feminine: \ar{حرب} \ar{طويلة}, a long war.

\begin{cousinweb}
\textbf{\ar{حرب}} (\emph{ḥarb}) is built on the native Arabic root \emph{ḥ-r-b} (Lane), where \emph{ḥaraba} is \textquotedblleft{}to strip a man of his goods\textquotedblright{}. War, which strips people of their goods, is usually explained from it.
\end{cousinweb}

\begin{grammarlens}[title={What you've built}]
One more word for telling what happened.
\end{grammarlens}

\subsection*{Your turn}
\begin{itemize}
\raggedright
  \item \emph{Say it:} \emph{ḥarb}
  \item \emph{Say it:} \emph{ḥarb}, once more
  \item \emph{Say it:} \emph{kāna jaddī ṣaghīran fī zaman al-ḥarb.} — My grandfather was a small boy in wartime.
  \item \emph{Recall:} say the Arabic for an army, then the Arabic for a soldier, then say \emph{ḥarb} again
\end{itemize}

\begin{tcolorbox}[breakable,colback=teal!4,colframe=teal!35!black,title={Before you move on}]
What does \ar{حرب} mean? (\textquotedblleft{}A war\textquotedblright{}.) Say it once more.
\end{tcolorbox}
\section[\texorpdfstring{\ar{معركة} (maʿraka)}{معركة (maʿraka)}]{\ar{معركة} (maʿraka) — a battle}

\label{lesson:AR-C307-maraka}



\begin{quote}
Before the new one: say the Arabic for a soldier, then the Arabic for a war.
\end{quote}

\subsection*{You'll want to know: \ar{معركة}}
\textbf{\ar{معركة}} — \emph{maʿraka} — \textquotedblleft{}a battle\textquotedblright{}.

\begin{quote}
\ar{وقعت} \ar{معركة} \ar{كبيرة} \ar{قرب} \ar{النهر}.
\end{quote}

\begin{quote}
\emph{waqaʿat maʿraka kabīra qurba an-nahr.}
\end{quote}

\begin{quote}
A big battle took place near the river.
\end{quote}

Plural: \ar{معارك}, \emph{maʿārik}.

\begin{cousinweb}
\textbf{\ar{معركة}} (\emph{maʿraka}) is built on the native Arabic root \emph{ʿ-r-k} (Lane), \textquotedblleft{}to rub hard, to grind\textquotedblright{}. \emph{Maʿraka} is the place where fighters grind against each other: the field of battle.
\end{cousinweb}

\begin{grammarlens}[title={What you've built}]
One more word for telling what happened.
\end{grammarlens}

\subsection*{Your turn}
\begin{itemize}
\raggedright
  \item \emph{Say it:} \emph{maʿraka}
  \item \emph{Say it:} \emph{maʿraka}, once more
  \item \emph{Say it:} \emph{waqaʿat maʿraka kabīra qurba an-nahr.} — A big battle took place near the river.
  \item \emph{Recall:} say the Arabic for a soldier, then the Arabic for a war, then say \emph{maʿraka} again
\end{itemize}

\begin{tcolorbox}[breakable,colback=teal!4,colframe=teal!35!black,title={Before you move on}]
What does \ar{معركة} mean? (\textquotedblleft{}A battle\textquotedblright{}.) Say it once more.
\end{tcolorbox}
\section[\texorpdfstring{\ar{حارب} (ḥāraba)}{حارب (ḥāraba)}]{\ar{حارب} (ḥāraba) — to wage war; to fight against}

\label{lesson:AR-C307-haraba-war}



\begin{quote}
Before the new one: say the Arabic for a war, then the Arabic for a battle.
\end{quote}

\subsection*{You'll want to know: \ar{حارب}}
\textbf{\ar{حارب}} — \emph{ḥāraba} — \textquotedblleft{}to wage war; to fight against\textquotedblright{}.

\begin{quote}
\ar{حارب} \ar{الأمير} \ar{جيشا} \ar{كبيرا}.
\end{quote}

\begin{quote}
\emph{ḥāraba al-amīr jayshan kabīran.}
\end{quote}

\begin{quote}
The prince fought a large army.
\end{quote}

Present: \ar{يحارب}, \emph{yuḥārib}.

It also fights a problem: \ar{حاربت} \ar{المدينة} \ar{التلوث}.

\begin{cousinweb}
\textbf{\ar{حارب}} (\emph{ḥāraba}) is built on the native Arabic root \emph{ḥ-r-b} (Lane), the root of \emph{ḥarb}, war. It is Form III, \textquotedblleft{}to make war with someone\textquotedblright{}.
\end{cousinweb}

\begin{grammarlens}[title={What you've built}]
One more word for telling what happened.
\end{grammarlens}

\subsection*{Your turn}
\begin{itemize}
\raggedright
  \item \emph{Say it:} \emph{ḥāraba}
  \item \emph{Say it:} \emph{ḥāraba}, once more
  \item \emph{Say it:} \emph{ḥāraba al-amīr jayshan kabīran.} — The prince fought a large army.
  \item \emph{Recall:} say the Arabic for a war, then the Arabic for a battle, then say \emph{ḥāraba} again
\end{itemize}

\begin{tcolorbox}[breakable,colback=teal!4,colframe=teal!35!black,title={Before you move on}]
What does \ar{حارب} mean? (\textquotedblleft{}To wage war; to fight against\textquotedblright{}.) Say it once more.
\end{tcolorbox}
